% =========================================================================
% Appendix B: Prompt Engineering for Automated Visual Complexity Scoring
% Two-column IEEE VIS / CHI style.
% Requires in the preamble:
%   \usepackage{booktabs}
%   \usepackage{array}
% Uses \cite{lin1989ccc}.
% Tables use tabular* with \extracolsep{\fill} so cells stay content-tight
% but inter-column gaps expand to fill the full column / page width.
% Narrow tables: \begin{table}  + tabular*{\linewidth}  (single column).
% Wide tables:   \begin{table*} + tabular*{\textwidth} (spans both columns).
% =========================================================================

\appendix
\section{Prompt Engineering for Visual Complexity Scoring}
\label{app:vc-prompts}

We investigate how the structure of a large multimodal LLM prompt affects the
model's ability to reproduce human-rated Visual Complexity (VC) scores for
static data visualizations. Two complementary experiments are reported:
\textbf{E1}, an ablation of four prompt components (Topics, Calibration,
Anchors, Weighted scoring), and \textbf{E2}, a comparison of four
qualitatively distinct prompt designs (V0--V3). Both experiments share the
same ground truth and evaluation protocol.

\subsection{Ground Truth and Evaluation Set}
\label{app:vc-gt}

The 46-image ground-truth set was sub-sampled from a pool of approximately
700 paired images collected in a prior pair-comparison study. Pairs were
ranked by the absolute difference $|\Delta \mathrm{VC}|$ between their two
aggregated human ratings, then binned across the full range of
$|\Delta \mathrm{VC}|$. From each bin we selected roughly equal numbers of
pairs, yielding \textbf{23 pairs (46 images)} that span the full difficulty
spectrum --- from near-ties (hard discrimination) to large-gap pairs (easy
discrimination). Each image carries a normalized human VC rating
$y \in [0,1]$ derived from the pair study.

Three of the 46 images are reserved as few-shot anchor exemplars
($\mathrm{VC} = 0.22, 0.54, 0.95$) and are presented in-context to every
anchored prompt (V2, V3, and the E1 variants containing ``A''). To prevent
leakage, these three images are \emph{excluded} from all reported metrics,
leaving $n = 43$ test images.

\subsection{Evaluation Metrics}
\label{app:vc-metrics}

We report five agreement metrics, with Lin's Concordance Correlation
Coefficient (CCC)~\cite{lin1989ccc} as the headline measure:
\begin{equation}
\mathrm{CCC}(y, \hat{y}) =
\frac{2\,\rho\,\sigma_y\sigma_{\hat{y}}}
     {\sigma_y^2 + \sigma_{\hat{y}}^2 + (\mu_y - \mu_{\hat{y}})^2},
\end{equation}
where $\rho$ is the Pearson correlation between human ratings $y$ and model
predictions $\hat{y}$. CCC measures agreement with the identity line
$\hat{y}=y$, jointly penalizing (i) low correlation, (ii) mean bias, and
(iii) scale mismatch. It is bounded in $[-1, 1]$ and is the standard metric
for continuous method-comparison studies. We additionally report Pearson $r$
with its two-sided $p$-value (\texttt{scipy.stats.pearsonr}), Spearman $\rho$,
and the coefficient of determination $R^2$ against the identity line
(\texttt{sklearn.metrics.r2\_score}),
$R^2 = 1 - \sum_i (y_i - \hat{y}_i)^2 / \sum_i (y_i - \bar{y})^2$, which can
be negative when the predictor underperforms the mean of $y$. MAE, RMSE, and
the signed bias $\overline{\hat{y}-y}$ complete the diagnostics.

\subsection{E1 --- Ablation of Prompt Components}
\label{app:vc-e1}

\subsubsection{Prompt components.}
Starting from a minimal \textbf{V0} baseline that states only the definition
of visual complexity and the 0--1 output range, we incrementally layer four
modular components. All E1 variants output only a single holistic
\texttt{vc\_score} (no per-dimension numeric scores).

\begin{itemize}
  \item \textbf{T (Topics / Dimensions):} Enumerates seven complexity
    dimensions (data density, visual encoding, semantics/text, schema,
    color/symbol, aesthetics, cognitive load) as mental checkpoints.
  \item \textbf{C (Calibration):} Provides scale-anchoring guidance with
    example chart types mapped to approximate target ranges
    (plain bar chart $\approx 0.25$--$0.40$; dense multi-panel
    $\approx 0.85$--$0.95$) and instructs the model to use the full range.
  \item \textbf{A (Anchors):} Prepends the three labeled anchor images as
    few-shot exemplars before the target. Each anchor carries only its final
    VC score (no per-dimension values).
  \item \textbf{W (Weighted):} Qualitatively re-weights the seven dimensions:
    high weight on density / encoding / schema / cognitive, medium on
    color / aesthetic, low on text.
\end{itemize}
All $2^3$ combinations of $(T, C, A)$ are tested, plus V0+TCA+W (9 variants).

\subsubsection{Test conditions.}
Each variant is evaluated under four independent runs on
\texttt{claude-opus-4-6} (primary model):
\texttt{r1} and \texttt{r2} are deterministic ($t=0$) replicates; \texttt{r3}
and \texttt{r4} use Claude's adaptive extended-thinking mode. A subset of
\texttt{r3} (variants V0+TC, V0+TA, V0+CA) used
\texttt{claude-sonnet-4-6}; these were re-run in \texttt{r4} on
\texttt{claude-opus-4-6} for cross-variant parity. We define
\emph{det\_avg} as the mean of \texttt{r1} and \texttt{r2}, and
\emph{think\_avg} as the mean of \texttt{r3} and \texttt{r4}.

\subsubsection{Results.}
Table~\ref{tab:e1-ccc} (column-width) summarizes CCC across all runs and
aggregated columns. Table~\ref{tab:e1-full} (page-width) reports the full
metric suite for \emph{det\_avg}.

% ---------- Narrow CCC table for E1 ----------------------------------------
\begin{table}[t]
  \centering
  \caption{E1 --- CCC per run and aggregated averages.
    \emph{det\_avg} averages \texttt{r1}+\texttt{r2} ($t=0$);
    \emph{think\_avg} averages \texttt{r3}+\texttt{r4} (thinking).
    Higher is better.}
  \label{tab:e1-ccc}
  \small
  \begin{tabular*}{\linewidth}{@{\extracolsep{\fill}}lcccccc@{}}
    \toprule
    Variant & r1 & r2 & \textbf{det\_avg} & r3 & r4 & \textbf{think\_avg}\\
    \midrule
    V0        & 0.858 & 0.859 & \textbf{0.858} & 0.868 & 0.881 & \textbf{0.875}\\
    V0+T      & 0.871 & 0.867 & \textbf{0.869} & 0.857 & ---   & \textbf{0.857}\\
    V0+C      & 0.851 & 0.856 & \textbf{0.854} & ---   & ---   & ---\\
    V0+A      & 0.852 & 0.853 & \textbf{0.852} & ---   & ---   & ---\\
    V0+TC     & 0.854 & ---   & \textbf{0.854} & 0.737 & 0.866 & \textbf{0.802}\\
    V0+TA     & 0.860 & ---   & \textbf{0.860} & 0.849 & 0.834 & \textbf{0.841}\\
    V0+CA     & 0.852 & ---   & \textbf{0.852} & 0.831 & 0.873 & \textbf{0.852}\\
    V0+TCA    & 0.856 & 0.856 & \textbf{0.856} & 0.874 & 0.873 & \textbf{0.874}\\
    V0+TCA+W  & 0.868 & 0.871 & \textbf{0.869} & 0.811 & 0.815 & \textbf{0.813}\\
    \bottomrule
  \end{tabular*}
\end{table}

% ---------- Wide full-metrics table for E1 ---------------------------------
\begin{table*}[t]
  \centering
  \caption{E1 --- Full metrics on \emph{det\_avg} (mean of two $t=0$
    replicates, \texttt{claude-opus-4-6}). $n=43$ non-anchor images.
    $R^2$ is \texttt{sklearn.metrics.r2\_score} against the identity line
    ($1 - \mathrm{SS}_\mathrm{res}/\mathrm{SS}_\mathrm{tot}$).
    All Pearson $p < 10^{-12}$.}
  \label{tab:e1-full}
  \small
  \begin{tabular*}{\textwidth}{@{\extracolsep{\fill}}lcccccccc@{}}
    \toprule
    Variant & $n$ & CCC & $r$ & $\rho$ (Spearman) & $R^2$ & MAE & RMSE & Bias\\
    \midrule
    V0        & 43 & 0.858 & 0.868 & 0.868 & 0.671 & 0.080 & 0.098 & $+0.010$\\
    V0+T      & 43 & 0.869 & 0.879 & 0.872 & 0.692 & 0.078 & 0.095 & $+0.003$\\
    V0+C      & 43 & 0.854 & 0.871 & 0.873 & 0.679 & 0.081 & 0.097 & $+0.033$\\
    V0+A      & 43 & 0.852 & 0.869 & 0.866 & 0.638 & 0.083 & 0.103 & $-0.012$\\
    V0+TC     & 43 & 0.854 & 0.874 & 0.870 & 0.688 & 0.078 & 0.096 & $+0.037$\\
    V0+TA     & 43 & 0.860 & 0.882 & 0.872 & 0.656 & 0.081 & 0.101 & $-0.025$\\
    V0+CA     & 43 & 0.852 & 0.863 & 0.859 & 0.653 & 0.084 & 0.101 & $+0.014$\\
    V0+TCA    & 43 & 0.856 & 0.863 & 0.866 & 0.674 & 0.079 & 0.098 & $+0.007$\\
    V0+TCA+W  & 43 & 0.869 & 0.876 & 0.871 & 0.704 & 0.075 & 0.093 & $-0.003$\\
    \bottomrule
  \end{tabular*}
\end{table*}

\subsubsection{Findings.}
At $t=0$, prompt engineering produces only marginal gains: all nine variants
land within a $\sim\!0.02$ CCC band ($[0.852, 0.869]$), with V0+T and
V0+TCA+W tied at the deterministic best ($0.869$) and the baseline V0
($0.858$) trailing by only $0.011$. Extended thinking helps minimal prompts
(V0: $0.858 \to 0.875$; V0+TCA: $0.856 \to 0.874$) but \emph{hurts}
V0+TCA+W ($0.869 \to 0.813$); the qualitative weight guidance appears to
interact poorly with longer deliberation. The \texttt{r3} sonnet runs
(V0+TC / V0+TA / V0+CA) are recovered by their \texttt{r4} opus replicates
(e.g., V0+TC: $0.737 \to 0.866$), confirming the drop is a model-capacity
effect, not a prompt defect.

\subsection{E2 --- Prompt-Version Comparison (V0--V3)}
\label{app:vc-e2}

\subsubsection{Prompt versions.}
E2 compares four qualitatively distinct prompt designs. Unlike E1, \textbf{V1
and V3 require the model to emit numeric scores for each of the seven
dimensions} in addition to the overall \texttt{vc\_score}:

\begin{itemize}
  \item \textbf{V0} --- Pure zero-shot: VC definition only, no dimensions,
    no anchors. (Identical to E1's V0 baseline.)
  \item \textbf{V1} --- Zero-shot with the seven dimension descriptions and
    detailed per-dimension scoring. No anchors, no calibration, no weighting.
  \item \textbf{V2} --- VC definition plus the three anchor images (anchors
    carry only overall VC, no per-dimension supervision).
  \item \textbf{V3} --- Full prompt: dimensions + calibration + weighting +
    three anchors with full per-dimension reference scores, and detailed
    per-dimension output.
\end{itemize}

The two principal differences between V3 and E1's V0+TCA+W are that
(a) V3's anchors include per-dimension ground-truth values, and (b) V3
requires the model to emit per-dimension scores --- neither of which is
present in V0+TCA+W.

\subsubsection{Test conditions.}
All four versions were run on \texttt{claude-opus-4-6} with the Anthropic
API default sampling temperature ($t=1$, not explicitly set) and no extended
thinking. \texttt{max\_tokens = 800} for V0 and V2;
\texttt{max\_tokens = 1500} for V1 and V3 to accommodate the longer
per-dimension JSON output. Each version is evaluated from a single
generation pass. Relative to E1, E2 matches the model but differs in
sampling (E1: $t=0$ with two replicates; E2: $t=1$, single run); E2 results
therefore carry more sampling noise than E1's \emph{det\_avg}.

\subsubsection{Results.}
Table~\ref{tab:e2-ccc} (column-width) shows the CCC summary;
Table~\ref{tab:e2-full} (page-width) reports the full metric suite.

% ---------- Narrow CCC table for E2 ----------------------------------------
\begin{table}[t]
  \centering
  \caption{E2 --- CCC for the four prompt versions.
    $n=43$ non-anchor images, \texttt{claude-opus-4-6}, $t=1$, single run.}
  \label{tab:e2-ccc}
  \small
  \begin{tabular*}{\linewidth}{@{\extracolsep{\fill}}lcc@{}}
    \toprule
    Version & $n$ & \textbf{CCC}\\
    \midrule
    V0          & 43 & 0.858\\
    V1          & 43 & 0.802\\
    V2          & 43 & 0.852\\
    \textbf{V3} & 43 & \textbf{0.882}\\
    \bottomrule
  \end{tabular*}
\end{table}

% ---------- Wide full-metrics table for E2 ---------------------------------
\begin{table*}[t]
  \centering
  \caption{E2 --- Full metrics on all four prompt versions.
    $n=43$ non-anchor images, \texttt{claude-opus-4-6}, $t=1$, single run.
    $R^2$ is \texttt{sklearn.metrics.r2\_score} against the identity line.}
  \label{tab:e2-full}
  \small
  \begin{tabular*}{\textwidth}{@{\extracolsep{\fill}}lccccccccc@{}}
    \toprule
    Version & $n$ & CCC & $r$ & $p$ (Pearson) & $\rho$ (Spearman) & $R^2$ & MAE & RMSE & Bias\\
    \midrule
    V0          & 43 & 0.858 & 0.867 & $5.5\!\times\!10^{-14}$ & 0.869 & 0.670 & 0.080 & 0.099 & $+0.011$\\
    V1          & 43 & 0.802 & 0.831 & $5.6\!\times\!10^{-12}$ & 0.816 & 0.613 & 0.086 & 0.107 & $-0.043$\\
    V2          & 43 & 0.852 & 0.868 & $4.8\!\times\!10^{-14}$ & 0.865 & 0.637 & 0.083 & 0.103 & $-0.012$\\
    \textbf{V3} & 43 & \textbf{0.882} & \textbf{0.900} & $\mathbf{2.0\!\times\!10^{-16}}$ & \textbf{0.896} & \textbf{0.718} & \textbf{0.076} & \textbf{0.091} & $-0.025$\\
    \bottomrule
  \end{tabular*}
\end{table*}

\subsubsection{Findings.}
V3 dominates every metric (CCC $0.882$, $r = 0.900$, $R^2 = 0.718$,
MAE $0.076$): $+0.024$ CCC over V0 and $+0.030$ over V2.
V1 --- per-dimension scoring without anchor exemplars --- is the
\emph{worst} version (CCC $0.802$, bias $-0.043$); requiring numeric
per-dimension output destabilizes calibration when no reference images are
provided. Anchors alone (V2, $0.852$) match the V0 baseline ($0.858$); the
benefit of anchors only emerges when paired with per-dimension supervision
(V3).

\subsection{Cross-Experiment Summary}
\label{app:vc-cross}

All runs use \texttt{claude-opus-4-6} (Table~\ref{tab:cross-exp}). Four
findings:
(1)~Prompt components alone have small, inconsistent effects
(E1: 0.02 CCC band at $t=0$).
(2)~Extended thinking is a double-edged sword: it helps minimal prompts
(V0, V0+TCA) but hurts heavily-instructed ones (V0+TCA+W).
(3)~The single change that reliably breaks the $0.87$ ceiling is requiring
per-dimension scoring \emph{with} per-dimension anchor supervision (V3).
Per-dimension output without anchors (V1) is actively harmful; anchors
without per-dimension output (V2) are neutral.
(4)~Because E1 and E2 share the same model, V3's advantage over every E1
variant is a clean prompt-design effect, not a model-capacity confound.

\begin{table}[t]
  \centering
  \caption{Cross-experiment summary. All runs use \texttt{claude-opus-4-6}.}
  \label{tab:cross-exp}
  \small
  \begin{tabular*}{\linewidth}{@{\extracolsep{\fill}}llccc@{}}
    \toprule
    Prompt & Exp. / Condition & CCC & $r$ & MAE\\
    \midrule
    V0 (baseline)        & E1 $t{=}0$ (det\_avg)    & 0.858 & 0.868 & 0.080\\
    V0+T                 & E1 $t{=}0$ (det\_avg)    & 0.869 & 0.879 & 0.078\\
    V0+TCA+W             & E1 $t{=}0$ (det\_avg)    & 0.869 & 0.876 & 0.075\\
    V0                   & E1 thinking (avg)        & 0.875 & 0.889 & 0.077\\
    V0+TCA+W             & E1 thinking (avg)        & 0.813 & 0.876 & 0.090\\
    V0                   & E2 $t{=}1$ (single run)  & 0.858 & 0.867 & 0.080\\
    V1                   & E2 $t{=}1$ (single run)  & 0.802 & 0.831 & 0.086\\
    V2                   & E2 $t{=}1$ (single run)  & 0.852 & 0.868 & 0.083\\
    \textbf{V3}          & \textbf{E2 $t{=}1$ (single run)} & \textbf{0.882} & \textbf{0.900} & \textbf{0.076}\\
    \bottomrule
  \end{tabular*}
\end{table}

======== BibTeX entry for bibliography file ===============================
@article{lin1989ccc,
  author  = {Lin, Lawrence I-Kuei},
  title   = {A Concordance Correlation Coefficient to Evaluate Reproducibility},
  journal = {Biometrics},
  volume  = {45},
  number  = {1},
  pages   = {255--268},
  year    = {1989},
  doi     = {10.2307/2532051},
  publisher = {International Biometric Society},
}
